\chapter{Linear Theory}
\label{ch:linear}

Chapters~5--8 developed the geometric and variational structure of
lamphrodynamic fields $(\Phi,\vec v,S)$, culminating in the classical
Hamiltonian formulation.
We now turn to analytic questions, beginning with the linearization of
the Euler--Lagrange equations \eqref{eq:EL-Phi}--\eqref{eq:EL-S} about
a smooth reference configuration and establishing global well--posedness
in Sobolev spaces.
This chapter collects the functional--analytic tools needed for the
nonlinear theory of Chapters~10--11.



\section{Linearization about a Reference Configuration}
\label{sec:linearization}

Let $(\Phi,\vec v,S)$ be a smooth lamphrodynamic configuration solving
the classical Euler--Lagrange equations on a globally hyperbolic
spacetime $(M,g)$.
Define the linearized fields by
\[
(\phi, w_{\mu}, s)
:=
\frac{d}{d\varepsilon}\big|_{\varepsilon=0}
(\Phi+\varepsilon\phi,\; v_{\mu}+\varepsilon w_{\mu},\; S+\varepsilon s).
\]
To simplify notation, write $U:=(\phi, w, s)$.
Substituting into \eqref{eq:EL-Phi}--\eqref{eq:EL-S} and retaining
terms linear in $U$ yields a linear system of wave--type operators
\begin{equation}
\label{eq:linear-system}
\mathcal{L}_{\mathrm{lin}} U
= \square_{g} U
 + B^{\mu}\nabla_{\mu}U
 + C\,U
=0,
\end{equation}
where $\square_{g}$ is the tensor wave operator, $B^{\mu}$ is a
first--order coefficient determined by the background
$(\Phi,\vec v,S)$, and $C$ is a zero--order coefficient.
The coefficients $B^{\mu},C$ are smooth and bounded on compact sets.



\begin{remark}[Hyperbolic Character]
The principal part of \eqref{eq:linear-system} is the wave operator,
ensuring hyperbolicity; lower--order terms encode entropy--flux
corrections of the background.
\end{remark}



\section{Functional Spaces and Initial Value Problem}
\label{sec:spaces}

Fix a spacelike Cauchy hypersurface~$\Sigma$ with induced Riemannian
metric~$h$.
For integer $k\ge0$ define Sobolev norms
\[
\|U\|_{H^{k}(\Sigma)}
 :=
 \sum_{j=0}^{k}
  \| \nabla^{j} U \|_{L^{2}(\Sigma)}.
\]

\begin{definition}[Initial Value Problem]
\label{def:IVP}
Given initial data
\[
U(0)=U_{0}\in H^{k}(\Sigma),\qquad
\partial_{t}U(0)=U_{1}\in H^{k-1}(\Sigma),
\]
find $U(t,x)$ solving \eqref{eq:linear-system} for $t$ in some time
interval containing $0$.
\end{definition}



\section{Energy Estimates}
\label{sec:linear-energy}

Define the linear energy functional
\begin{equation}
\label{eq:linear-energy}
E_{\mathrm{lin}}(U)(t)
:= \frac{1}{2}
\int_{\Sigma_{t}}
\bigl(
 |\partial_{t}U|^{2}
 +|\nabla U|^{2}
 + |U|^{2}
\bigr)\, d\mathrm{vol}_{h}.
\end{equation}
Standard integration--by--parts yields:

\begin{lemma}[Linear Energy Estimate]
\label{lem:linear-energy}
For sufficiently smooth $U$ solving \eqref{eq:linear-system},
\[
\frac{d}{dt}E_{\mathrm{lin}}(U)(t)
\le
C\bigl(E_{\mathrm{lin}}(U)(t)+\|U(t)\|_{L^{2}}^{2}\bigr),
\]
where $C$ depends on the background coefficients.
\end{lemma}

\begin{proof}[Proof (Sketch)]
Multiply \eqref{eq:linear-system} by $\partial_{t}U$,
integrate over $\Sigma_{t}$,
use divergence theorem on the wave operator, and absorb
lower--order terms into the right--hand side.
\end{proof}

By Grönwall's inequality we obtain:

\begin{corollary}[A priori Control]
\label{cor:apriori}
For $t\in[0,T]$,
\[
E_{\mathrm{lin}}(U)(t)
\le
E_{\mathrm{lin}}(U)(0)\,\exp(C T).
\]
\end{corollary}



\section{Global Well--posedness}
\label{sec:global-wp}

We now state the main linear result.

\begin{theorem}[Global Well--posedness in Sobolev Spaces]
\label{thm:linear-wp}
Let $(M,g)$ be globally hyperbolic and let $k\ge2$.
For any initial data $(U_{0},U_{1})\in H^{k}(\Sigma)\times H^{k-1}(\Sigma)$
there exists a unique global solution
\[
U\in C^{0}([0,\infty);H^{k}(\Sigma))
 \cap C^{1}([0,\infty);H^{k-1}(\Sigma))
\]
of \eqref{eq:linear-system}.
Moreover, the energy bound of Corollary~\ref{cor:apriori} holds.
\end{theorem}

\begin{proof}[Proof]
Local existence and uniqueness follow from standard hyperbolic PDE
theory for systems with smooth coefficients
\cite{Hormander-AnalysisPDE}.
Global extension follows from the a priori bound
(Corollary~\ref{cor:apriori}) combined with continuation criteria.
\end{proof}



\section{Spectral Properties of the Linearized Operator}
\label{sec:spectral}

On static spacetimes or after Fourier transform in time,
$\mathcal{L}_{\mathrm{lin}}$ reduces to an elliptic operator of the
form
\[
\mathscr{A} = -\Delta_{h} + V(x),
\]
where $V$ is a potential induced by the background.
Standard spectral theory implies:

\begin{proposition}[Spectral Properties]
\label{prop:spectral}
If $(M,g)$ is static and $\Sigma$ compact, then
$\mathscr{A}$ has discrete spectrum $\{\lambda_{j}\}_{j\ge0}$,
with $\lambda_{j}\to+\infty$ as $j\to\infty$.
\end{proposition}

\begin{proof}
Immediate from ellipticity and compactness of $\Sigma$
\cite{Taylor-PDE}.
\end{proof}



\section{Connection to AKSZ Classical Limit}
\label{sec:connection-AKSZ-classical}

The linear theory developed here governs fluctuations around classical
solutions of the AKSZ--RSVP action. The coercivity of the Hamiltonian
(Chapter~8) and the linear estimate above ensure that the BV
perturbative expansion is formally well behaved around smooth
backgrounds, a prerequisite for quantization in Volume~III.



\section{Summary}

We have shown that linearized lamphrodynamic fields satisfy a
hyperbolic system of second--order PDEs with global well--posedness
in Sobolev spaces. The resulting energy and spectral control forms the
analytic basis for the nonlinear theory of Chapters~10--11.

